% SOURCE_AUTHORITY: sga5_fr_workpass.tex, lines 5347--5375 (LF-normalized unit).
% SOURCE_SHA256: 791F4EFFC5E02832D5D77ED03518C8156D6F07E4C8238B03545DB93D883FBB28
\subsubsection*{5.11. Aplicación a los grupos finitos: propiedades de divisibilidad}

Sea $G$ un grupo finito, cuyo elemento neutro se denota por $e$, y pongamos $K[G]=B$. El $K$-módulo $B_\natural$ es libre con base el conjunto $G_\natural$ de las clases de conjugación de $G$; toda sección $s$ de $B_\natural$ se escribe de manera única
\[
s=\sum_{x\in G_\natural}s(x)x.
\]

\begin{statement}{Proposición 5.11.1}
Cf. SGA $4^{1/2}$, Rapport 4.5. Sean $E\in\ob D(B)_{\mathrm{parf}}$, $M\in\ob D^b(K)$. Para $u\in\uHom_B(E,M\lten_KE)$, se tiene, en $H^0(T,M)$,
\begin{equation}\tag{*}
\Tr_K(u)=|G|\footnote{Si $X$ es un conjunto, se denota por $|X|$ el cardinal de $X$.}\,\Tr_B(u)(e).
\end{equation}
\end{statement}

En efecto, basta aplicar 5.10.11 con $A=K$, observando que $\Tr_{B/K}:B_\natural\to K$ viene dada por
\begin{equation}\tag{5.11.2}
\Tr_{B/K}(g)=
\begin{cases}
|G|,& g=e,\\
0,& g\ne e.
\end{cases}
\end{equation}

\begin{statement}{Proposición 5.11.3}
Bajo las hipótesis de 5.11.1, sea $g\in G$, denotemos por $Z_g$ el centralizador de $g$ en $G$ y pongamos $A=K[Z_g]$. Sea $u\in\uHom_B(E,M\lten_KE)$; denotemos por $gu\in\uHom_A(E,M\lten_KE)$ la compuesta de la imagen de $u$ en $\uHom_A(E,M\lten_KE)$ con la multiplicación a la izquierda por $g$. Entonces se tiene, en $H^0(T,M)$,
\[
\Tr_A(gu)(e)=\Tr_B(u)(g^{-1}).
\]
\end{statement}
